\chapter{The language and its semantics}
\label{chap:semantics}

% Use one small constructor-and-alternative example throughout this chapter.
% Present the mathematical definitions before discussing their Coq encoding.

\section{Expressions, patterns, and matchings}
% Introduce the calculus and show how a Haskell definition becomes a matching.
% Explain the normalized syntax required for sharing during lazy evaluation.

\section{Arity and weak-head normal forms}
% Define matching arity, partial application, and weak-head normal forms.
% Distinguish expression values from matching success and failure.

\section{Big-step evaluation}
% Present the mutually dependent expression and matching judgments.
% Explain heap updates, black-holing, and left-to-right alternatives through
% the running example.

\section{The small-step machine}
% Introduce configurations, the two control modes, and continuation frames.
% Follow the same example through machine transitions and compare it with the
% big-step derivation. Leave Coq representation choices to the next chapter.
